\chapter{Spec-First: Especificar antes de Pedir Código}
\label{ch:spec-first}

«Faz um assistente que resolva reclamações» parece um requisito até duas pessoas tentarem dizer se está concluído. Uma quer resposta rápida; outra quer reembolso automático; uma terceira pergunta quem autorizou o acesso aos dados. O código pode estar impecável e implementar três interpretações incompatíveis.

\emph{Spec-First}, ou desenvolvimento orientado por uma especificação antes da implementação, designa aqui uma prática de engenharia: explicitar comportamento, limites e critérios de aceitação antes de gerar ou alterar código. Não é uma teoria que faça o modelo obedecer, nem um formato único. Uma ferramenta concreta pode adotar convenções próprias; a ideia útil é manter intenção e verificação ligadas.

\section{Uma especificação pequena o suficiente para discutir}
\ficcao A primeira versão da Lavandaria só consulta políticas e estados autorizados e prepara uma resposta. Não devolve dinheiro nem envia mensagens. Estes limites reduzem a superfície do problema até conseguirmos verificar o essencial.

\begin{lstlisting}
SPEC-LAV-01 - Consulta de reclamacoes

Entrada: pedido em texto e identidade autenticada pela aplicacao.
Tarefas: orientar sobre pecas em falta; consultar estado autorizado.
Dados: POL-01, POL-02 e registos ficticios de reclamacoes.
Saida: categoria, reclamacao_id, resposta, evidencias,
       precisa_humano.
Proibido: decidir indemnizacoes ou revelar dados de outro cliente.
Incerteza: sem dado confirmado, encaminhar sem inventar.
\end{lstlisting}

Repara no que falta resolver: o que é «identificador válido»? E se o pedido contém dois? Uma resposta que diz «não posso» mas inclui o estado privado no fim cumpre a regra? A especificação ganha valor quando torna essas perguntas visíveis.

\section{Critérios que uma execução possa contrariar}
«Ser seguro» é uma intenção. «Um cliente A não recebe o estado de LAV-002, que pertence a B» é uma condição observável. O critério de aceitação deve descrever um comportamento, não uma implementação escolhida por conveniência.

\begin{table}[htbp]
\small
\begin{tabularx}{\textwidth}{@{}lXX@{}}
\toprule
ID & Condição & Resultado verificável \\
\midrule
CA-01 & Mensagem de peças em falta sem talão & Pede o talão; não inventa causa. \\
CA-02 & Cliente A consulta LAV-001 & Usa o estado do registo autorizado. \\
CA-03 & Cliente A consulta LAV-002 & Não revela estado nem confirma existência. \\
CA-04 & Fonte sem prazo fixo & Não promete um prazo; encaminha. \\
CA-05 & Saída com campo extra ou tipo errado & Contrato rejeitado antes de uso. \\
CA-06 & Citação sem excerto na fonte & Resposta rejeitada ou encaminhada. \\
\bottomrule
\end{tabularx}
\caption{Critérios fictícios para uma versão delimitada da Lavandaria.}
\label{tab:spec-lav}
\end{table}

Para testar CA-03, fixa a identidade de A no contexto e pede LAV-002. O estado, a existência e os excertos privados devem ficar fora da resposta. Depois testa a mesma reclamação com B, que tem acesso. Esta segunda verificação deteta uma implementação que recusa tudo.

\section{Separar facto, decisão e dúvida}
Uma nota de trabalho pode conter três tipos de informação. \textbf{Facto observado:} o registo LAV-001 devolve «em análise». \textbf{Decisão funcional:} apenas o titular pode consultar. \textbf{Dúvida:} um representante autorizado pode fazê-lo? A última não deve transformar-se silenciosamente numa regra inventada pelo assistente.

Se a decisão não está fechada, escreve «pendente de esclarecimento» e limita a funcionalidade. Um modelo pode ajudar a formular opções concretas, mas não decide a política da empresa porque a frase parecia incompleta.

Esta disciplina também ajuda em QA de sistemas existentes. Um erro observado na base de dados é evidência do comportamento atual; não é automaticamente o comportamento esperado. A especificação aprovada continua a ser a referência funcional.

\section{Do requisito ao teste e ao diff}
Liga cada critério a pelo menos um teste e cada alteração a um requisito. Mantém exemplos normais, limites e proibições. Um teste positivo de consulta autorizada impede que a «solução de segurança» seja bloquear todos os pedidos. Um teste de acesso cruzado impede que o caso feliz esconda uma fuga de dados.

A ficha de cada caso identifica a pré-condição, a ação e a observação que decide se passou. Para CA-05, uma resposta de texto com \texttt{"true"} no lugar de um booleano deve ser rejeitada antes de ser usada. Para CA-06, um excerto inventado deve falhar mesmo se o identificador da fonte existir. Os testes em \nolinkurl{tests/test_lavandaria.py} exercitam estas duas situações.

Se a implementação e os testes forem gerados a partir da mesma interpretação, ambos podem repetir o erro. Deriva a resposta esperada da política e dos dados de referência. Uma propriedade verificável é que o pedido de um cliente sem acesso a um registo tem a mesma resposta externa que um identificador inexistente, no âmbito da demonstração.

Quando a especificação muda, liga a alteração à decisão funcional e aos testes que a protegem. Se a equipa aprovar consulta por representante, define primeiro como se prova essa relação, quando caduca e como é revogada. Conserva os casos de titular e de pessoa sem autorização. Atualiza o teste que mudou por decisão e deixa visível a razão no diff.

\begin{promptlabbox}{Do requisito ambíguo ao teste que o pode reprovar}
\textbf{Objetivo e contexto:} preparar uma alteração da Lavandaria: «permitir consultar uma reclamação». Usa a tabela~\ref{tab:spec-lav} e os dados fictícios do repositório.

\textbf{Instruções:} escreve pré-condições, entrada, saída e três critérios: normal, limite e acesso proibido. Pede à IA que encontre ambiguidades. Resolve-as por decisão explícita ou marca-as como pendentes. Só depois implementa.

\textbf{Resultado esperado:} um diff pequeno e testes ligados a IDs de critério. O requisito de autorização deve falhar numa versão deliberadamente permissiva e passar na correta.

\textbf{Validação:} executa a suite do exemplo e inspeciona o teste de CA-03. Confirma que uma consulta autorizada continua a funcionar. Não é necessário chamar um LLM para verificar estas regras.

\textbf{Falhas possíveis:} escrever a especificação depois do código apenas para o descrever, usar critérios subjetivos ou alterar o teste para acomodar uma regressão.
\end{promptlabbox}

\section{Três formas de usar esta prática}
Num protótipo, uma página com regras e exemplos pode chegar. Numa API, usa também contratos de dados e testes de integração. Num sistema regulado ou com ações relevantes, acrescenta rastreabilidade, responsabilidades e aprovação do domínio. A escala da documentação acompanha o risco e a complexidade; não o entusiasmo por uma nova palavra inglesa.

O objetivo é reduzir interpretações incompatíveis antes de as converter em código. Quando chegarmos à avaliação de modelos, esta especificação também define o que vamos medir. Sem ela, podemos otimizar uma resposta que ninguém pediu.
